\documentclass[10pt,english]{beamer}
%\documentclass[english,handout]{beamer} % For handouts
\usetheme[progressbar=frametitle,block=fill]{metropolis} %numbering=none

%%% USEFUL PACKAGES
%\usepackage{showframe} % For debugging positioning
\usepackage{etex} % If too many packages
% Encoding and language
\usepackage[utf8]{inputenc}
\usepackage{babel}
\usepackage{amsmath, amssymb}
\usepackage{natbib}
%\usepackage{booktabs}
%\usepackage{algorithmic}
\usepackage{algorithm}
\usepackage{caption}
%\usepackage{animate} % Animations
\usepackage{bm} % Bold math
\usepackage{bbm}
%\usepackage{url}
%\usepackage{pifont}
%\usepackage{ulem} % Used for strikeouts \sout
%\usepackage{stackengine}
%\usepackage{enumitem}
%\setlist[description]{leftmargin=\parindent,labelindent=\parindent}
%\usepackage{colortbl} % Used for colored rows in tables


%%% GRAPHICS
\usepackage{graphicx}
\graphicspath{{./figs/}}


%%% COLORS
\setbeamercolor{background canvas}{bg=white}
\def\BlankFrame{
	\bgroup
	%\pdfpageheight 29.7cm
	\setbeamercolor{background canvas}{bg=}
	\begin{frame}[plain]
	\end{frame}
	%\makeatletter
	%\pdfpageheight \beamer@paperheight
	%\makeatother
	\egroup}

\usepackage{xcolor}
\definecolor{DarkGreen}{HTML}{00B200}
\definecolor{LightBlue}{HTML}{0090D9}
\definecolor{gold}{rgb}{.812,.710,.231}
% Text markup
%\setbeamercolor{alerted text}{fg=red}
\newcommand{\blue}[1]{\textcolor{blue}{#1}}
\newcommand{\red}[1]{\textcolor{red}{#1}}
\newcommand{\grey}[1]{\textcolor{gray}{#1}}
\newcommand{\orange}[1]{\textcolor{mLightBrown}{#1}}
\newcommand\myheading[1]{\textbf{#1}}
\newcommand\myemph[1]{\underline{\emph{#1}}}
\newcommand\textexample[1]{\textit{\textbf{#1}}}

%%% SPACING
\newcommand\vws[1][1]{\vspace{#1\baselineskip}} % vertical white space
%\newcommand\strt[1][1.5ex]{\rule[-.05\baselineskip]{0pt}{#1}} % strut
\newcommand\strt[2]{\rule[-#1ex]{0pt}{#2ex}} % strut
\newcommand\Hrule{\vspace{1ex} \hrule \vspace{1ex}} % Horisontal rule with some space after

%%% MISC
\newcommand\articleref[4]{\noindent\begin{minipage}[t]{0.04\textwidth}
		\vspace{0pt} 
		\pgfuseimage{beamericonarticle}
	\end{minipage}%
	\begin{minipage}[t]{0.96\textwidth}
		\vspace{0pt}
		#1. \textbf{#2.} \textit{#3}, #4.
	\end{minipage}}

%%% METROPOLIS THEME SPECIFIC
\makeatletter
\setlength{\metropolis@progressonsectionpage@linewidth}{1pt}
\makeatother
%\setbeamercolor{progress bar}{fg=red,bg=red!50}


%%% TEXTPOS
\usepackage[absolute,overlay]{textpos} % option showboxes is useful in draft mode
\setlength{\TPHorizModule}{\paperwidth}
\setlength{\TPVertModule}{\paperheight}
\textblockorigin{0pt}{10mm} % start everything at top-left, below gray 


%%% TIKZ/PGFPLOTS
\usepackage{tikz}
\usetikzlibrary{arrows,positioning,calc,shapes.geometric}
%\usetikzlibrary{arrows,calc,shapes.geometric,decorations.pathmorphing,backgrounds,positioning,fit,petri,decorations.pathreplacing}
%\usepackage{pgfplots}
%\pgfplotsset{compat = 1.3}


%%% BLOCKS AND BOXES
% Changing colors of blocks
%\setbeamercolor{block title alerted}{bg=UURed,fg=palette primary.fg}
%\setbeamercolor{block body alerted}{bg=UURed!15}
\setbeamercolor{block title alerted}{bg=mLightBrown,fg=palette primary.fg}
\setbeamercolor{block body alerted}{bg=mLightBrown!15}
%\setbeamercolor{block title example}{bg=UUGreen,fg=palette primary.fg}
%\setbeamercolor{block body example}{bg=UUGreen!10}
% \mybox is a rectangular box
\usepackage{boxedminipage}
\setlength\fboxrule{2pt}
\setlength\fboxsep{2\fboxsep}
\newcommand\mybox[3][\textwidth]{
  {\color{#2}
    \begin{boxedminipage}{#1}
      {\color{palette primary.bg} #3}
    \end{boxedminipage}}%
}   
\usepackage{tcolorbox}
\tcbset{arc=1mm,grow to left by=3mm,grow to right by=3mm,left=2mm}
%\newenvironment{redbox}{%
%	\begin{tcolorbox}[colback=UURed!15,colframe=UURed]}{%
%	\end{tcolorbox}}
%\newenvironment{greenbox}{%
%	\begin{tcolorbox}[colback=UUGreen!15,colframe=UUGreen]}{%
%	\end{tcolorbox}}
\newenvironment{redbox}{%
	\begin{tcolorbox}[colback=red!15,colframe=red]}{%
	\end{tcolorbox}}
\newenvironment{greenbox}{%
	\begin{tcolorbox}[colback=DarkGreen!15,colframe=DarkGreen]}{%
	\end{tcolorbox}}
\newenvironment{graybox}{%
	\begin{tcolorbox}[colback=mDarkTeal!5,colframe=mDarkTeal]}{%
	\end{tcolorbox}}
\newenvironment{orangebox}{%
\begin{tcolorbox}[colback=mLightBrown!15,colframe=mLightBrown]}{%
	\end{tcolorbox}}
\newenvironment{bwbox}{%
	\begin{tcolorbox}[colback=white,colframe=black]}{%
\end{tcolorbox}}
\newenvironment{bluebox}{%
	\begin{tcolorbox}[colback=LightBlue!15,colframe=LightBlue]}{%
\end{tcolorbox}}


%%%%%%%%% NEW MACROS

\newcommand\imp[1]{\alert{\textbf{#1}}}
\newcommand\bfit[1]{\textbf{\textit{#1}}}
\newcommand\good{\color{DarkGreen}{$\blacktriangle$}} % used in lists
\newcommand\bad{\color{red}{$\blacktriangledown$}} % used in lists


\RequirePackage{amsmath, amssymb}
\RequirePackage{bbm}
%\RequirePackage{newtxmath}


% Convenience macro for referring to data source
\newcommand\sourceurl[2]{\small \grey{Data from \href{#1}{#2}}}

% Abbreviations
\RequirePackage{xspace}
\newcommand\pdf{pdf\xspace}
\newcommand\ifft{iff\xspace}
\newcommand\ex{\textbf{ex)}\xspace}

% General time series notation
\newcommand\T{n}  % Length of time series
\newcommand\rtheta{{\red{\theta}}}  % Parameter (color coded)
\newcommand\rthetah{{\red{\widehat\theta}}}  % Estimate (color coded)

% Neural netowkrs
\newcommand\h{\mathbf{h}} % Hidden state variable
\newcommand\zz{\mathbf{z}} % Generic input (vector)

% For OLS/AR
\newcommand\noise{\varepsilon}  % This is the noise in AR, but should it be the same as measurement noise in SSM?
\newcommand\noisevar{\sigma^2_\noise}
\newcommand\noisevarhat{\widehat\sigma^2_\noise}
\newcommand\X{\Phi}
\newcommand\y{\mathbf{y}}
\newcommand\bphi{\bm\phi}

% State space models
\newcommand\z{\alpha}  % State vector, general SSM
\newcommand{\obsnoise}{\varepsilon}
\newcommand{\statenoise}{\eta}
\newcommand{\varobs}{\sigma^2_{\varepsilon}}
\newcommand{\varstate}{\sigma^2_{\eta}}
% For structural time series
\newcommand{\trendnoise}{\zeta}
\newcommand{\seasnoise}{\omega}
\newcommand{\vartrend}{\sigma^2_{\trendnoise}}
\newcommand{\varseas}{\sigma^2_{\seasnoise}}

%
\newcommand\FF{T}
\newcommand\GG{R}
\newcommand\HH{Z}
\newcommand{\covobs}{\sigma_\epsilon^2}
\newcommand{\covstate}{Q}
\newcommand\initmean{a_1}
\newcommand\initcov{P_1}
% Kalman filter
\newcommand{\zpart}[2]{\z_{#1}^{#2}}
\newcommand{\wgt}[2]{\omega_{#1}^{#2}}
\newcommand{\wgtsum}[1]{\Omega_{#1}}
\newcommand\zhat[2]{\hat\z_{#1|#2}}
\newcommand\Phat[2]{P_{#1|#2}}
\newcommand\zpred[1]{\zhat{#1}{#1-1}}
\newcommand\Ppred[1]{\Phat{#1}{#1-1}}
\newcommand\zfilt[1]{\zhat{#1}{#1}}
\newcommand\Pfilt[1]{\Phat{#1}{#1}}
\newcommand\ypred[1]{\hat y_{#1|#1-1}}
\newcommand\Spred[1]{F_{#1|#1-1}}
\newcommand\Spredinv[1]{\Spred{#1}^{-1}}
\newcommand\epshat[2]{\hat{\obsnoise}_{#1|#2}}
\newcommand\etahat[2]{\hat{\statenoise}_{#1|#2}}

\newcommand{\statefun}{T}
\newcommand{\obsfun}{Z}
\newcommand{\estfun}{h}

\newcommand{\qd}{q} %State density
\newcommand{\md}{g} %Measure density

\newcommand{\rmd}{\mathrm{d}}

% SMC
\newcommand{\Np}{N}           % Number of particles
\newcommand{\Mp}{M}           % Number of particles in backward simulation



%\RequirePackage{color}
%\newcommand{\flnote}[1]{{\color{red}\textbf{[#1]}}} % Used for notes in text - color red
%\newcommand\Hrule{\vspace{1ex} \hrule \vspace{1ex}} % Horisontal rule with some space after; This is moved to beamer preamble

%%%%%%%%%%%%%%%%%%%%%%%%%%%%%%%%%%%%%%%%%%%%%%%%%%%%%%%%%%%%%%%%%%%%%%%%%%%%%%%%
%                            COMMANDS IN TEXT                                  %
%%%%%%%%%%%%%%%%%%%%%%%%%%%%%%%%%%%%%%%%%%%%%%%%%%%%%%%%%%%%%%%%%%%%%%%%%%%%%%%%
\newcommand\numtext[2]{#1\textsuperscript{#2}}
\newcommand\thsnd[1]{\ensuremath{#1\thinspace000}}
\newcommand{\peqref}[1]{\eqref{#1} on page~\pageref{#1}} % Page referencing for equations: "(1) on page 1"

%%%%%%%%%%%%%%%%%%%%%%%%%%%%%%%%%%%%%%%%%%%%%%%%%%%%%%%%%%%%%%%%%%%%%%%%%%%%%%%%
%                            SPECIFIC MATH                                     %
%%%%%%%%%%%%%%%%%%%%%%%%%%%%%%%%%%%%%%%%%%%%%%%%%%%%%%%%%%%%%%%%%%%%%%%%%%%%%%%%
% Models etc.
%\newcommand{\T}{T}            % Number of samples in data record
\newcommand{\parspace}{\Theta}                                   % Parameter space
\newcommand{\parameter}{\theta}                                  % Parameter
% Spaces
\newcommand{\setX}{\ensuremath{\mathsf{X}}}                      % State-space X
\newcommand{\sigmaX}{\ensuremath{\mathcal{X}}}                   % Sigma algebra on X
\newcommand{\setY}{\ensuremath{\mathsf{Y}}}                      % State-space Y
\newcommand{\sigmaY}{\ensuremath{\mathcal{Y}}}                   % Sigma algebra on Y
\newcommand{\setZ}{\ensuremath{\mathsf{Z}}}                      % State-space Z
\newcommand{\sigmaZ}{\ensuremath{\mathcal{Z}}}                   % Sigma algebra on Z

%%%%%%%%%%%%%%%%%%%%%%%%%%%%%%%%%%%%%%%%%%%%%%%%%%%%%%%%%%%%%%%%%%%%%%%%%%%%%%%%
%                           GENERAL MATH                                       %
%%%%%%%%%%%%%%%%%%%%%%%%%%%%%%%%%%%%%%%%%%%%%%%%%%%%%%%%%%%%%%%%%%%%%%%%%%%%%%%%

% ======== Miscellaneous symbols ========
\newcommand\eqdef{:=}
\newcommand\defeq{=:}
\newcommand\const{\text{const.}}
%\newcommand\eqdef{\stackrel{\text{\scriptsize def}}{=}}

\newcommand\iid{iid}
\newcommand{\iidsim}{\stackrel{\text{\iid}}{\sim}} % iid simulation
\newcommand{\process}[1]{\{#1\}_{t\geq 1}}       % Process (time index t)
\newcommand{\range}[2]{#1, \, \dots, \, #2}      % Range = 1, ..., N
\newcommand{\crange}[2]{\{#1, \, \dots, \, #2\}} % Curly range = {1, ..., N}
\newcommand{\prange}[2]{(#1, \, \dots, \, #2)}   % Parenthesised range = (1, ..., N)
\newcommand{\bwdrange}[2]{#1 : -1 : #2}          % Range = N, ..., 1
\newcommand{\approxpropto}{\stackrel{\sim}\propto}

% Tight dots between \int and \int in a multidimensional integral
\newcommand{\tightcdots}{\hspace*{-0.38em}\cdot\hspace*{-0.3em}\cdot\hspace*{-0.3em}\cdot\hspace*{-0.38em}}

% Arrows - convergence and mappings
% \mapsto                                                     % Mappings, x \mapsto f(x)
\newcommand{\fromto}{\rightarrow}                             % Mapping from set A to set B; f: A \fromto B
\newcommand{\goesto}{\rightarrow}                             % limits used in n \goesto \infty
\newcommand{\goestosmall}{\to}                                % limits used in \lim_{n \goestosmall \infty}
\newcommand{\convP}{\stackrel{\probab}\longrightarrow}        % Convergence in probability
\newcommand{\convD}{\stackrel{\textrm{D}}\longrightarrow}     % Convergence in distribution

% ======== Standard spaces  ========
\newcommand{\naturals}{\ensuremath{\mathbb{N}}}               % Natural numbers
\newcommand{\reals}{\ensuremath{\mathbb{R}}}                  % Real numbers
\newcommand{\nonnegatives}{\reals_{\smaller +}}               % Nonnegative numbers
\newcommand{\positives}{\reals_{\smaller ++}}                 % Positive numbers
\newcommand{\nonnegativedefinites}[1]{S_{\smaller +}(#1)}     % Nonnegative #1 x #1 matrices
\newcommand{\positivedefinites}[1]{S_{++}(#1)}                % Positive #1 x #1 matrices

% ======== Matrices ========
\newcommand{\eye}[1]{I_{#1}}                     % Identity matrix
\newcommand{\+}{\mathsf{T}}                      % Transpose
\newcommand{\kronecker}{\raisebox{1pt}{\ensuremath{\otimes}}} % Kronecker product
\DeclareMathOperator*\diag{diag}
\DeclareMathOperator*\trace{tr}

% ======== Operators, calculus etc. ========
\newcommand{\Ordo}{O}                            % Big ordo
\newcommand{\supnorm}[1]{\|#1\|_\infty}          % Supremum norm
\newcommand\osc{\text{osc}}                      % Oscillator norm
\newcommand{\grad}{\nabla}                       % Gradient
\newcommand{\complementof}[1]{\ensuremath{#1^\mathsf{c}}} % Set complement
\renewcommand\vec{\text{vec}}
\DeclareMathOperator*\supp{supp}                          % Support
\DeclareMathOperator*\card{card}                          % Set cardinality
\DeclareMathOperator*\rank{rank}                          % Rank
\DeclareMathOperator*\sign{sign}                          % Signum function
\DeclareMathOperator*\argmax{arg\,max}
\DeclareMathOperator*\argmin{arg\,min}

% ======== Probability ========
\newcommand{\Prb}{\ensuremath{\mathbb{P}}}                       % Probability
\newcommand{\E}{\ensuremath{\mathbb{E}}}                         % Expectation
\newcommand{\var}{\ensuremath{\mathrm{Var}}}                     % Variance
\newcommand{\cov}{\ensuremath{\mathrm{Cov}}}                     % Covariance
\newcommand{\cor}{\ensuremath{\mathrm{Corr}}}                     % Correlation
\newcommand{\I}{\ensuremath{\mathbbm{1}}}						 % Indicator function

%\newcommand{\abscont}{\ensuremath{\ll}}          % Absolute continuity
\renewcommand\mid{\,\vert\,} % I don't really like that \mid produces rubber lengths. Sometimes, we get very large white spaces p(x    |   y), and it can produce line breaks after "p(x |" . Is the non-rubber definition here better?
\newcommand\Mid{\,\middle\vert\,} % Stretchable |, to use with \left \right - N.B. This produces a longer | in general. Does that look better than a standard \mid?


% Distributions
\newcommand{\N}{\ensuremath{\mathcal{N}}}        % Normal
\newcommand{\uni}{\ensuremath{\mathcal{U}}}      % Uniform
\newcommand\MN{\mathcal{MN}}                     % Matrix normal
\newcommand\IW{\mathcal{IW}}                     % Inverse-Wishart
\newcommand\GP{\mathcal{GP}}                     % Gaussian process
\DeclareMathOperator*\Mult{Mult}                 % Multinomial
\DeclareMathOperator*\cat{Cat}                   % Categorical
\DeclareMathOperator*\Discrete{Discrete}         % Categorical/alternative name
\DeclareMathOperator*\bin{Bin}                   % Binomial
\DeclareMathOperator*\gam{Gam}                   % Gamma
\DeclareMathOperator*\St{St}                     % Student's t
\DeclareMathOperator*\po{Po}                   % Binomial

%\usepackage{extendedalt}
%\usepackage{animate} % Animations
%\usepackage{../lindsten}
%\usepackage{movie15}
\usepackage{tikz}
\usepackage{listofitems} % for \readlist to create arrays

\title{732G12 Data Mining}
\subtitle{Föreläsning 5}
\date{}
\author{Johan Alenlöv \\ IDA, Linköping University, Sweden}
\titlegraphic{\hfill\includegraphics[height=1.2cm]{../LiU_primary_black.pdf}}
%\institute{Joint work with\dots}


%% MY DEF %%
\newcommand{\itm}[1]{\mathrm{Item}_{#1}}
\newcommand{\pausa}{\pause}
%\renewcommand{\pausa}{}
\tikzstyle{mynode}=[thick,draw=blue,fill=blue!20,circle,minimum size=22]


\newenvironment{nscenter}
 {\parskip=0pt\par\nopagebreak\centering}
 {\par\noindent\ignorespacesafterend}

\begin{document}

\maketitle

\begin{frame}{Dagens föreläsning}

    \begin{itemize}
        \item Neurala nätverk
        \item Feature learning
        \item Optimering av neurala nätverk
        \item Hyperparametrar
    \end{itemize}
    
\end{frame}

\begin{frame}{Tentainformation}
    
    Finns några gamla tentor på Lisam. Den tenta ni får kommer vara liknande men inget som har med associationsanalys/sekventiell data är med i kursen nu.

    Inte bara kod som ska lämnas in utan även lösningar. Använd Rmarkdown för att skriva fina lösningar med plottar och kod. All kod ska bifogas i inlämningen.

    På Lisam finns mapp med alla hjälpfiler för tentan. Dessa kommer finnas på datorerna när tentan börjar. Kolla igenom dessa i god tid innan tentan för att få en uppfattning om vad som finns och var!

\end{frame}

\begin{frame}{Neuralt nätverk}
    När vi pratar om neurala nätverk kan vi prata om lite vad som helst. Finns \textbf{väldigt många} olika sorters nätverk.

    Kan använda neurala nätverk för att lösa många olika problem:
    \begin{itemize}
        \item Övervakad inlärning
        \item Oövervakad inlärning
        \item Reinforcement learning
        \item Generativa modeller
        \item Representation learning
    \end{itemize}

\end{frame}


\begin{frame}{Neurala nätverk}
    
    För övervakad inlärning kan vi t.ex. använda
    \begin{itemize}
        \item Feed-forward network / Multiple layer perceptron (MLP)
        \item Radial basis network
        \item Convolutional neural networks (CNN)
        \item Recurrent neural networks
    \end{itemize}

\end{frame}


\begin{frame}{Neurala Nätverk}

    För oövervakad inlärning:
    \begin{itemize}
        \item Dolda representationer: Autoencoders
        \item Clustering: Self Organizing Map
    \end{itemize}

    Generativa modeller:
    \begin{itemize}
        \item Används för att lära sig komplexa fördelningar för att sen dra nya samples.
        \begin{itemize}
            \item Sampla nya bilder
            \item Skriva text
        \end{itemize}
        \item Generative adversarial network (GAN)
    \end{itemize}
    
\end{frame}

\begin{frame}{Feature learning}

    Vi går tillbaka till linjär regression,
    \begin{equation*}
        y = \mathbf{X} \beta + \varepsilon, \quad \mathbb{E}[\varepsilon] = 0, \quad \mathbb{V}[\varepsilon] = \sigma^2.
    \end{equation*}
    
    Vad kan vi göra om data inte följer denna linjära modell?

\end{frame}

\begin{frame}{Feature learning}
    
    Vanligt i linjär regression,
    \begin{itemize}
        \item Givet data $\mathbf{X} = (x_1, x_2, \ldots, x_p)$ och $ y = \mathbf{X} \beta$.
        \item Vi kan transformera variablerna i $\mathbf{X}$:
        \begin{itemize}
            \item Polynomregression, $\mathbf{X} = (x, x^2, x^3, \ldots, x^p)$
            \item Funktioner, $\log(x), \sqrt{x}, \exp(x), \ldots$
            \item Interaktioner, $x_1 x_2$
            \item Stegfunktioner
            \item Disktretisering
            \item Dummy-kodning
        \end{itemize}
        \item Kallas i maskininlärning för "feature engineering"
        \begin{itemize}
            \item Svårt att veta vilken transformation som vi ska göra för varje problem.
            \item Svårt med komplexa datastrukturer som text eller bild.
        \end{itemize}
    \end{itemize}

\end{frame}

\begin{frame}{Feature learning}
    
    \begin{itemize}
        \item Vi har data $\mathbb{X} = (x_1, x_2, \ldots, x_p)$.
        \item Transformationer är funktioner av data.
        \begin{itemize}
            \item Ex. $h(x) = \log(x), h(x_1,x_2) = \exp(x_1 \cdot x_2)$.
        \end{itemize}
        \item Anta en $x$-variabel, låt $h(x)$ vara en viktad summa av andra funktioner,
        \begin{equation*}
            z = h(x) = \sum_{i=1}^{M}w_i h_i(x),
        \end{equation*}
        där $h_i(x)$ är godtyckliga funktioner.
        \item<2-> Om vi har många $x$-variabler får vi,
        \begin{equation*}
            z =  h(x_1, x_2, \ldots, x_p) = \sum_{i=1}^{M} w_i h_i(x_1, x_2, \ldots, x_p).
        \end{equation*}
        \item<2-> Hur ska vi välja $h_i(x)$?
    \end{itemize}

\end{frame}

\begin{frame}{Feature learning}
    
    För en linjär transformation hitta matriserna $\mathbf{W}$ och $\mathbf{V}$,
    \begin{equation*}
        \underset{n \times m}{\mathbf{Z}} = \underset{n \times p}{\mathbf{X}} \cdots \underset{p \times m}{\mathbf{W}}, \qquad \underset{n \times g}{\mathbf{Z}} = \underset{n \times p}{\mathbf{X}} \cdots \underset{p \times m}{\mathbf{W}} \cdot \underset{m \times g}{\mathbf{V}}.
    \end{equation*}

    För neurala nätverk vill vi kunna modellera icke-linjära funktioner.

    Idé: Använd många "enkla" icke-linjära funktioner för att skapa en komplex icke-linjär funktion!

\end{frame}

\begin{frame}{Feature learning}

    \begin{greenbox}
        \myheading{Neurala nätverk}

        Låt $\sigma(\cdot)$ vara en enkel icke-linjär funktion och låt $h_i(x_1, \ldots, x_p)$ vara en linjär funktion,
        \begin{equation*}
            h_i(x_1, x_2, \ldots, x_p) = \beta_{0i} + \beta_i^{\top} \mathbf{x}.
        \end{equation*}

        Låt nu
        \begin{equation*}
            z = \sigma(h_i(x_1, x_2, \ldots, x_p)) = \sigma(\beta_{0i} + \beta_i^{\top} \mathbf{x}),
        \end{equation*}
        nästla sedan många sådana funktioner för att bygga upp en godtyckligt komplex icke-linjär funktion.
    \end{greenbox}
    
\end{frame}

\begin{frame}{Feature learning}
    
    För MLP brukar vi skriva
    \begin{equation*}
        \underset{k \times 1}{a}^{(p+1)} = \sigma \left( \underset{k \times n}{\mathbf{W}}^{(p)} \cdot \underset{n \times 1}{a}^{(p)} + \underset{k \times 1}{b}^{(p)} \right).
    \end{equation*}
    Här är:
    \begin{description}
        \item[$k$] Dimension av nya lagret
        \item[$n$] Dimension av föregående lager
        \item[$\mathbf{W}$] Viktmatris
        \item[$b$] Bias
        \item[$(p)$] Vilket lager
        \item[$\sigma()$] Vår funktion som opererar elementvis   
    \end{description}
    Historiskt har sigmoid eller hyperbolic tangent varit vanliga aktiveringsfunktioner. Numera är ReLu (eller varianter) den vanligaste,
    \begin{equation*}
        \operatorname{ReLu}(x) = \max(0,x).
    \end{equation*}
\end{frame}

\begin{frame}{Feature learninig}

    Vi kan se ett neuralt nätverk som att vi
    \begin{enumerate}
        \item Automatiskt lär oss transformationer av de förklarande variablerna.
        \item Gör linjär (logistisk,multinomiell) regression på transformationerna (sista lagret).
    \end{enumerate}
    OBS!
    \begin{itemize}
        \item Komplexa funktioner kräver mycket data att lära sig!
        \item Neurala nätverk kan lätt överanpassa träningsdata!
        \item Funkar när vi har stort antal förklarande variabler.
        \item Om vi låter gömda lager ha mindre dimension än förklarande variabler får vi "icke-linjär variabelreduktion".
    \end{itemize}
    
\end{frame}

\begin{frame}{Universal approximation theorem}
    Vilka funktioner kan vi då lära oss med ett neuralt nätverk av detta slag?

    \begin{greenbox}
        \myheading{Universal approximation theorem}: Let $\operatorname{C}(\mathrm{X},\mathbb{R}^{m})$ denote the set of continuous functions from $\mathrm{X} \subset \mathbb{R}^n$ to $\mathbb{R}^{m}$. Let $\sigma \in \operatorname{C}(\mathbb{R},\mathbb{R})$. Note that $(\sigma \otimes x)_i = \sigma(x_i)$, so $\sigma \otimes x$ denotes $\sigma$ applied to each component of $x$.

        Then $\sigma$ is not polynomial if and only if for every $n \in \mathbb{N}, m \in \mathbb{N}$, compact $\mathrm{K} \subseteq \mathbb{R}^n, f \in \operatorname{C}(\mathrm{K},\mathbb{R}^m), \varepsilon > 0$ there exists $k \in \mathbb{N}, A \in \mathbb{R}^{k \times n}, b \in \mathbb{R}^k, C \in \mathbb{R}^{m \times k}$ such that
        \begin{equation*}
            \sup_{x \in \mathrm{K}} \| f(x) - g(x) \| < \varepsilon,
        \end{equation*}
        where $g(x) = C \cdot (\sigma \otimes (A \cdot x + b))$.
    \end{greenbox}
\end{frame}

\begin{frame}{Optimering av neurala nätverk}
    
    Gradient decent: Hitta minimum på en funktion genom att gå dit den lutar mest!

    Vill hitta
    \begin{equation*}
        a^* = \argmin_a L(a) = \sum_i L_i(f(x^{(i)},a), y^{(i)}).
    \end{equation*}

    Löser detta genom sekvensen
    \begin{equation*}
        a_{n+1} = a_n - \gamma \cdot \nabla L(a_n).
    \end{equation*}

    \begin{itemize}
        \item Vi behöver gradienter (partiella derivator)
        \item Backpropagation: kedjeregel för derivator på neruala nätverk.
        \item Gradient decent: dyrt när vi har många observationer!
    \end{itemize}

\end{frame}

\begin{frame}{Optimering av neurala nätverk}
    
    Svårt problem med många fallgropar.

    \begin{itemize}
        \item Lokala minima
        \begin{itemize}
            \item Ställen som ser ut som ett minima (eller grannar har högre kostand) men inte är det bästa som finns.
            \item Kan ha hög kostand eller låg.
            \item Identifikationsproblem:
            \begin{itemize}
                \item Viktsymmetri
                \item Skalning mellan lager
            \end{itemize}
            \item Kan ha oräkneligt antal lokala minima.
        \end{itemize}
    \end{itemize}

\end{frame}

\begin{frame}{Optimering av neurala nätverk}
    
    Platåer och sadelpunkter
    \begin{itemize}
        \item Ställen där gradienten är noll (eller nära), fast vi inte är på ett loaklt min/max.
        \item Sadelpunkter:
        \begin{itemize}
            \item Lokalt minima i några riktningar.
            \item Lokalt maxima i andra riktningar.
        \end{itemize}
        \item Antalet sadelpunkter tenderar att öka med antalet dimensioner.
        \item Platåer är stora områden som är platta (gradient nära noll).
        \item Platåer och sadelpunkter gör optimeringen med gradient decent svårare.
    \end{itemize}

\end{frame}

\begin{frame}{Stochastic gradient decent (SGD)}

    Det är dyrt att beräkna $\nabla L(a_n)$ för alla datapunkter.

    Gör istället en väntesvärdesriktig skattning $\nabla \hat{L}(a_n)$ av gradienten genom att ta ett slumpmässigt sample från data (mini-batch).

    \begin{itemize}
        \item Större batch ger mindre varians i skattningen men blir dyrare att beräkna.
        \item Kräver fler iterationer och mindre learning rate.
        \item Kräver att vi har oberoende observationer.
        \item Funkar bra för neurala nätverk!
        \item En epoch är en genomgång av all träningsdata.
    \end{itemize}
    
\end{frame}

\begin{frame}{Hyperparametrar}
    
    I neruala nätverk finns massvis med Hyperparametrar!
    \begin{itemize}
        \item Arkitektur:
        \begin{itemize}
            \item Antal gömda lager
            \item Antal neuroner i varje lager
            \item Aktiveringsfunktioner
            \item (Speciella typer av neuroner/lager)
        \end{itemize}
        \item Optimeringen:
        \begin{itemize}
            \item Mini-batch storlek
            \item Learning rate (fix eller föränderlig)
            \item Antal epoker
            \item (Vilken optimeringsalgoritm som används)
        \end{itemize}
    \end{itemize}

\end{frame}

\begin{frame}{Hyperparametrar}
    Hur ska vi bestämma deras värden?
    \begin{itemize}
        \item Svår fråga utan exakt svar.
        \item Mycket trial and error.
        \item Valideringsdata
        \item För stora problem kan det ta lång tid att hitta bra hyperparametrar
    \end{itemize}
\end{frame}

\end{document}